%lezione del 04/03/2026
%tenuta dal prof costa del corso A: non capisco come unire le lezioni in un discorso
%coerente e sensato: per ora le butto insieme e vedrò
\newpage
NOTE TO SELF: qui iniziano le lezioni di Costa (04/03), non c'è un filo che lega a prima,
da rivedere per cercare di darci un senso, probabilmente farà K-G dopo.
\section{Teoria degli urti}
Nella fisica nucleare e subnucleare, gli urti giocano un ruolo fondamentale:
a differenza della meccanica classica, descrivono dei processi in maniera molto
precisa (quando due biglie collidono non sono davvero puntiformi, quando lo fanno due
elettroni sì): per trattarli si userà la cinematica relativistica.

Si adotta da ora in poi, se non specificato, la notazione per cui il quadrimomento di una
particella $ a $  viene denotato con lo stesso nome della particella:
$ a = (\frac{E_a}{c},\vec{a}) $.

Si consideri ora l'urto tra due particelle $ a $ e $ b $; esse possono urtare in due modi: 
\begin{itemize}
    \item Urto elastico: le particelle rimangono le medesime, ma i loro quadrimomenti cambiano;	
    \item Urto anelastico: le particelle possono cambiare a seguito dell'interazione.
\end{itemize}
In entrambi i casi il quadrimomento totale è conservato:
$$ a + b = a' + b',\quad a + b = e + f + g $$ 
nel caso, rispettivamente, di urto elastico e non elastico.

Si noti che la quantità $ a^2 = \frac{E_a^2}{c^2} - |\vec{a}| = m_a^2c^2$ è un invariante,
quindi è una caratteristica della particella, detta \textbf{massa invariante}.

\subsection{Cinematica degli urti elastici}
Si considera ora il semplice caso di una particella $ a $ che urta contro un target, costituito
da particelle $ b $: la conservazione del quadrimomento è
$$ a + b = a' + b' $$
Vale anche la conservazione dei moduli quadri
$$ (a+b)^2 = (a'+b')^2 \rightarrow a^2 + b^2 +2a\cdot b =a'^2 + b'^2 +2a'\cdot b' $$
poiché $ a $ e $ a' $ sono la stessa particella, la loro massa invariante è uguale:

$a^2 = a'^2$, similmente $b^2 = b'^2$.
